\section*{Inhaltsangabe.}

\subsection*{Wesen und Begriff der Theologie.}
\emph{(De Natura et Constitutione Theologiae.)}

\begin{enumerate}
    \item Die Verständigung über den Standpunkt, S. 1. -- 2. Über Religion im allgemeinen, S. 6. -- 3. Die Zahl der Religionen in der Welt, S. 8. -- 4. Die zwei Erkenntnisquellen der tatsächlich bestehenden Religionen, S. 10. -- 5. Die Ursache der Parteien in der äußeren Christenheit, S. 23. -- 6. Das Christentum als absolute Religion, S. 36. -- 7. Christliche Religion und christliche Theologie, S. 42. -- 8. Die christliche Theologie, S. 44. -- 9. Die nähere Beschreibung der Theologie, als Tüchtigkeit gefasst, S. 50. -- 10. Die nähere Beschreibung der Theologie, als Lehre gefasst, S. 56. -- 11. Einteilungen der Theologie, als Lehre gefasst, S. 84. Gesetz und Evangelium, S. 84. Fundamentale und nichtfundamentale Lehren, S. 89. Primäre und sekundäre Fundamentallehren, S. 95. Nichtfundamentale Lehren, S. 102. Offene Fragen und theologische Probleme, S. 104. -- 12. Die Kirche und die kirchlichen Dogmen, S. 108. -- 13. Der Zweck der Theologie, den sie an dem Menschen erreichen will, S. 116. -- 14. Die äußeren Mittel der Theologie, wodurch sie ihr Ziel an den Menschen erreicht, S. 118. -- 15. Theologie und Wissenschaft, S. 119. -- 16. Theologie und Gewissheit, S. 123. -- 17. Theologie und Lehrfortbildung, S. 147. -- 18. Theologie und Lehrfreiheit, S. 154. -- 19. Theologie und System, S. 158. -- 20. Theologie und Methode, S. 172. -- 21. Die Erlangung der theologischen Tüchtigkeit, S. 228.
\end{enumerate}

\subsection*{Die heilige Schrift. (De Scriptura Sacra.)}

\begin{enumerate}
    \item Die heilige Schrift ist für die Kirche unserer Zeit die einzige Quelle und Norm der christlichen Lehre, S. 233. -- 2. Die heilige Schrift ist im Unterschiede von allen andern Schriften Gottes Wort, S. 256. -- 3. Die heilige Schrift ist Gottes Wort, weil sie von Gott eingegeben oder inspiriert ist, S. 262. -- 4. Das Verhältnis des Heiligen Geistes zu den Schreibern der heiligen Schrift, S. 275. -- 5. Die Einwände gegen die Inspiration der heiligen Schrift, S. 280 (der verschiedene Stil in den einzelnen Büchern der Schrift; die Berufung auf historische Forschung; die verschiedenen Lesarten; angebliche Widersprüche und irrige Angaben; ungenaue Zitate der neutestamentlichen Schreiber aus dem Alten Testament; geringe und dem Heiligen Geist nicht anständige Dinge; Solözismen, Barbarismen, verfehlte Satzkonstruktionen). -- 6. Geschichtliches zur Lehre von der Inspiration, S. 320. -- 7. Luther und die Inspiration der Schrift, S. 334. -- 8. Zusammenfassende Charakteristik der neueren Theologie, sofern sie die Inspiration der Schrift leugnet, S. 360. -- 9. Die Folgen der Leugnung der Inspiration, S. 367. -- 10. Die Eigenschaften der heiligen Schrift, S. 371 (die göttliche Auto-
\end{enumerate}